\section{Noetherian rings}
\label{section-Noetherian}

\noindent
A ring $R$ is {\it Noetherian} if any ideal of $R$ is finitely generated.
This is equivalent to the ascending chain condition for ideals of $R$.
Indeed, the union of an ascending sequence of ideals is an ideal; if
it has a finite generating set, all those generators lie in one term,
and the sequence stabilizes there. Conversely, if an ideal is not
finitely generated, start with $(0)$ and successively choose an element
of the ideal outside the ideal generated by all the preceding choices.
This gives a strictly increasing sequence of ideals, contradicting
the ascending chain condition.
By
Lemma \ref{lemma-cohen}
it suffices to check that every prime ideal of $R$ is finitely generated.

\begin{lemma}
\label{lemma-Noetherian-permanence}
\begin{slogan}
Noetherian property is stable by passage to finite type extension
and localization.
\end{slogan}
Any finitely generated ring over a Noetherian ring
is Noetherian. Any localization of a Noetherian ring
is Noetherian.
\end{lemma}

\begin{proof}
The statement on localizations follows from the fact
that any ideal $J \subset S^{-1}R$ is of the form
$I \cdot S^{-1}R$. Any quotient $R/I$ of a Noetherian
ring $R$ is Noetherian because any ideal $\overline{J} \subset R/I$
is of the form $J/I$ for some ideal $I \subset J \subset R$.
Thus it suffices to show that if $R$ is Noetherian so
is $R[X]$. Suppose $J_1\subset J_2\subset\cdots$ is an ascending
chain of ideals in $R[X]$. For $i\geq1$ and $d\geq0$ define
$$
I_{i,d}=\{a\in R\mid aX^d+\sum_{e=0}^{d-1}b_eX^e\in J_i
\text{ for some }b_0,\ldots,b_{d-1}\in R\}.
$$
This contains zero, and addition and scalar multiplication of the
displayed polynomials show that it is an ideal. For a nonzero $a$,
the displayed polynomial has degree exactly $d$ and leading
coefficient $a$. Multiplication by $X^{d'-d}$ and the inclusion
$J_i\subset J_{i'}$ prove $I_{i,d}\subset I_{i',d'}$ when
$i\leq i'$ and $d\leq d'$.

\medskip\noindent
Put $K_d=\bigcup_{i\geq1}I_{i,d}$, an ideal because the union is
ascending in $i$. The sequence $K_0\subset K_1\subset\cdots$
stabilizes at some $d_0$. Each $K_d$ for $0\leq d\leq d_0$ has a
finite generating set. Every generator lies in some $I_{i,d}$;
taking the largest of these finitely many indices gives $i_d$ with
$I_{i_d,d}=K_d$ (take $i_d=1$ if the generating set is empty).
Let $i_0=\max_{0\leq d\leq d_0}i_d$. For $d\leq d_0$ and
$i\geq i_0$ we have $I_{i,d}=I_{i_0,d}=K_d$.
For $d>d_0$ the inclusions
$$
K_{d_0}=I_{i_0,d_0}\subset I_{i_0,d}\subset I_{i,d}
\subset K_d=K_{d_0}
$$
give the same equality. Thus the required $i_0$ works simultaneously
for every degree.

\medskip\noindent
For $i\geq i_0$ and $f\in J_i$, prove $f\in J_{i_0}$ by induction
on its degree, with $f=0$ treated separately. If $f\ne0$ has degree
$d$ and leading coefficient $a\ne0$, the equality
$I_{i,d}=I_{i_0,d}$ supplies $g\in J_{i_0}$ of degree $d$ with
the same coefficient $a$. Then $f-g\in J_i$ is zero or has degree
less than $d$. For $d=0$ it is zero; otherwise induction applies.
Hence $f-g\in J_{i_0}$ and $f\in J_{i_0}$. All $J_i$ with
$i\geq i_0$ are therefore equal, as required.
\end{proof}

\begin{lemma}
\label{lemma-Noetherian-power-series}
If $R$ is a Noetherian ring, then so is the formal power
series ring $R[[x_1, \ldots, x_n]]$.
\end{lemma}

\begin{proof}
The isomorphism
$R[[x_1,\ldots,x_{n+1}]]\cong R[[x_1,\ldots,x_n]][[x_{n+1}]]$
sends a coefficient family $(a_{\alpha,e})$ to
$\sum_{e\geq0}(\sum_{\alpha\in\mathbf Z_{\geq0}^n}a_{\alpha,e}
x_1^{\alpha_1}\cdots x_n^{\alpha_n})x_{n+1}^e$.
Its inverse reads off the same coefficients. For each fixed
multi-index, the Cauchy product has only finitely many contributing
pairs, so these mutually inverse maps preserve addition, multiplication
and the identity. It therefore suffices to prove that $R[[x]]$ is
Noetherian when $R$ is Noetherian. Let $I\subset R[[x]]$ be an ideal.
We have to show that $I$ is a finitely generated ideal.
For each nonnegative integer
$d$ denote $I_d = \{a \in R \mid ax^d + \text{h.o.t.} \in I\}$.
Then we see that $I_0 \subset I_1 \subset \ldots$ stabilizes as $R$
is Noetherian. Choose $d_0$ such that $I_{d_0} = I_{d_0 + 1} = \ldots$.
For each $d \leq d_0$ choose elements $f_{d, j} \in I \cap (x^d)$,
$j = 1, \ldots, n_d$ such that if we write
$f_{d, j} = a_{d, j}x^d + \text{h.o.t}$ then $I_d = (a_{d, j})$.
Denote $I' = (\{f_{d, j}\}_{d = 0, \ldots, d_0, j = 1, \ldots, n_d})$.
Then it is clear that $I' \subset I$. Pick $f \in I$.
Choose the coefficients $c_{d,i}$ successively for $d=0,\ldots,d_0$.
Before the step for $d$, the current remainder lies in $I\cap(x^d)$.
Its coefficient of $x^d$ belongs to $I_d$, so express it as
$\sum_{i=1}^{n_d}c_{d,i}a_{d,i}$ and subtract
$\sum_{i=1}^{n_d}c_{d,i}f_{d,i}$. The new remainder lies in
$I\cap(x^{d+1})$. Thus
$$
h_0=f-\sum_{d=0}^{d_0}\sum_{i=1}^{n_d}c_{d,i}f_{d,i}
\in I\cap(x^{d_0+1}).
$$
For every $e\geq1$, having constructed
$h_{e-1}\in I\cap(x^{d_0+e})$, its coefficient of $x^{d_0+e}$
lies in $I_{d_0+e}=I_{d_0}$. Choose $c_{i,e}\in R$ such that this
coefficient is $\sum_{i=1}^{n_{d_0}}c_{i,e}a_{d_0,i}$, and put
$$
h_e=h_{e-1}-\sum_{i=1}^{n_{d_0}}c_{i,e}x^e f_{d_0,i}
\in I\cap(x^{d_0+e+1}).
$$
These choices define actual formal power series
$C_i(x)=\sum_{e=1}^\infty c_{i,e}x^e\in R[[x]]$.
For any fixed degree $N$, the coefficient of $x^N$ in
$C_i(x)f_{d_0,i}$ receives contributions only from
$1\leq e\leq N-d_0$, a finite set. The remainder formula with
$d_0+e+1>N$ therefore proves, coefficient by coefficient, the identity
$$
f=\sum_{d=0}^{d_0}\sum_{i=1}^{n_d}c_{d,i}f_{d,i}
 +\sum_{i=1}^{n_{d_0}}\left(\sum_{e=1}^\infty c_{i,e}x^e\right)f_{d_0,i}.
$$
This is a finite $R[[x]]$-linear combination of the original chosen
generators, so $f$ belongs to $I'$. No assertion that $I'$ is closed
under limits is needed: the equality is an equality of all formal
coefficients with finitely many generator terms.
\end{proof}

\noindent
The following lemma, although easy, is useful because
finite type $\mathbf{Z}$-algebras come up quite often in
a technique called ``absolute Noetherian reduction''.

\begin{lemma}
\label{lemma-obvious-Noetherian}
Any finite type algebra over a field is Noetherian.
Any finite type algebra over $\mathbf{Z}$ is Noetherian.
\end{lemma}

\begin{proof}
This is immediate from Lemma \ref{lemma-Noetherian-permanence}
and the fact that fields are Noetherian rings and that
$\mathbf{Z}$ is a Noetherian ring (because it is a
principal ideal domain).
\end{proof}

\begin{lemma}
\label{lemma-Noetherian-finite-type-is-finite-presentation}
Let $R$ be a Noetherian ring.
\begin{enumerate}
\item Any finite $R$-module is of finite presentation.
\item Any submodule of a finite $R$-module is finite.
\item Any finite type $R$-algebra is of finite presentation over $R$.
\end{enumerate}
\end{lemma}

\begin{proof}
Let $M$ be a finite $R$-module. By
Lemma \ref{lemma-trivial-filter-finite-module}
we can find a finite filtration of $M$ whose successive quotients are
of the form $R/I$. Since any ideal is finitely generated, each of
the quotients $R/I$ is finitely presented. Hence $M$ is finitely
presented by
Lemma \ref{lemma-extension}.
This proves (1).

\medskip\noindent
Let $N \subset M$ be a submodule. As $M$ is finite, the quotient
$M/N$ is finite. Thus $M/N$ is of finite presentation by part (1).
Thus we see that $N$ is finite by Lemma \ref{lemma-extension} part (5).
This proves part (2).

\medskip\noindent
To see (3) note that any ideal of
$R[x_1, \ldots, x_n]$ is finitely generated by
Lemma \ref{lemma-Noetherian-permanence}.
\end{proof}

\begin{lemma}
\label{lemma-Noetherian-topology}
If $R$ is a Noetherian ring then $\Spec(R)$
is a Noetherian topological space, see Topology,
Definition \ref{topology-definition-noetherian}.
\end{lemma}

\begin{proof}
This is because any closed subset of $\Spec(R)$
is uniquely of the form $V(I)$ with $I$ a radical ideal,
see Lemma \ref{lemma-Zariski-topology}.
And this correspondence is inclusion reversing.
Thus the result follows from the definitions.
\end{proof}

\begin{lemma}
\label{lemma-Noetherian-irreducible-components}
\begin{slogan}
A Noetherian affine scheme has finitely many generic points.
\end{slogan}
If $R$ is a Noetherian ring then $\Spec(R)$
has finitely many irreducible components. In other words
$R$ has finitely many minimal primes.
\end{lemma}

\begin{proof}
By Lemma \ref{lemma-Noetherian-topology} and
Topology, Lemma \ref{topology-lemma-Noetherian}
we see there are finitely many irreducible components.
By Lemma \ref{lemma-irreducible} these correspond to
minimal primes of $R$.
\end{proof}

\begin{lemma}
\label{lemma-Noetherian-base-change-finite-type}
Let $R \to S$ be a ring map. Let $R \to R'$ be of finite type.
If $S$ is Noetherian, then the base change $S' = R' \otimes_R S$
is Noetherian.
\end{lemma}

\begin{proof}
By Lemma \ref{lemma-base-change-finiteness} finite type is stable under
base change. Thus $S \to S'$ is of finite type. Since $S$ is Noetherian we
can apply Lemma \ref{lemma-Noetherian-permanence}.
\end{proof}

\begin{lemma}
\label{lemma-Noetherian-field-extension}
Let $k$ be a field and let $R$ be a Noetherian $k$-algebra.
If $K/k$ is a finitely generated field extension then
$K \otimes_k R$ is Noetherian.
\end{lemma}

\begin{proof}
Since $K/k$ is a finitely generated field extension, there exists
a finitely generated $k$-algebra $B \subset K$ such that $K$ is
the fraction field of $B$. In other words, $K = S^{-1}B$
with $S = B \setminus \{0\}$. Then $K \otimes_k R = S^{-1}(B \otimes_k R)$.
Then $B \otimes_k R$ is Noetherian by
Lemma \ref{lemma-Noetherian-base-change-finite-type}.
Finally, $K \otimes_k R = S^{-1}(B \otimes_k R)$ is Noetherian by
Lemma \ref{lemma-Noetherian-permanence}.
\end{proof}

\noindent
Here are some fun lemmas that are sometimes useful.

\begin{lemma}
\label{lemma-subring-of-local-ring}
Let $R$ be a ring and $\mathfrak p \subset R$ be a prime.
There exists an $f \in R$, $f \not \in \mathfrak p$ such
that $R_f \to R_\mathfrak p$ is injective in each of the
following cases
\begin{enumerate}
\item $R$ is a domain,
\item $R$ is Noetherian, or
\item $R$ is reduced and has finitely many minimal primes.
\end{enumerate}
In (2), it suffices more generally that
$\ker(R\to R_{\mathfrak p})$ is finitely generated.
\end{lemma}

\begin{proof}
If $R$ is a domain, then $R \subset R_\mathfrak p$, hence $f = 1$ works.
More generally, suppose the kernel $I=\ker(R\to R_{\mathfrak p})$
is generated by $a_1,\ldots,a_t$. For each $a_j$, the localization
criterion gives $s_j\notin\mathfrak p$ with $s_ja_j=0$.
Set $f=\prod_{j=1}^ts_j$, using $f=1$ when $t=0$.
Then $f\notin\mathfrak p$ and $fI=0$.
If $r/f^e$ maps to zero in $R_{\mathfrak p}$, the image of $r$
is zero, so $r\in I$ and $fr=0$. Hence $r/f^e=0$ in $R_f$.
This proves injectivity, and applies in particular when $R$ is Noetherian.

\medskip\noindent
For (3), retain all the minimal primes
$\mathfrak p_1,\ldots,\mathfrak p_n$. For each $i$ with
$\mathfrak p_i\not\subset\mathfrak p$, choose
$f_i\in\mathfrak p_i\setminus\mathfrak p$, and put
$$
f=\prod_{\mathfrak p_i\not\subset\mathfrak p}f_i.
$$
The empty product is $1$, and primality of $\mathfrak p$ gives
$f\notin\mathfrak p$. If $r\in\ker(R\to R_{\mathfrak p})$,
some $s\notin\mathfrak p$ satisfies $sr=0$. Whenever
$\mathfrak p_i\subset\mathfrak p$, we have $s\notin\mathfrak p_i$,
so $r\in\mathfrak p_i$ by primality. For every other $i$ the
chosen factor $f_i$ gives $f\in\mathfrak p_i$.
It follows that $fr$ belongs to every minimal prime. The intersection
of the minimal primes is the nilradical, which is zero since $R$
is reduced. Thus $fr=0$. The same fraction argument as above proves
that the original canonical map $R_f\to R_{\mathfrak p}$ is injective.
\end{proof}

\begin{lemma}
\label{lemma-surjective-endo-noetherian-ring-is-iso}
Any surjective endomorphism of a Noetherian ring is an isomorphism.
\end{lemma}

\begin{proof}
Suppose $f:R\to R$ is surjective and not injective, and choose
$0\ne a\in\ker(f)$. For each $n\geq1$, surjectivity of $f^n$
provides $b_n\in R$ with $f^n(b_n)=a$. Then
$f^{n+1}(b_n)=f(a)=0$ but $f^n(b_n)\ne0$.
Thus every inclusion in
$$
\ker(f)\subset\ker(f^2)\subset\ker(f^3)\subset\cdots
$$
is strict, witnessed by $b_n$ at the $n$th inclusion. This contradicts
the ascending chain condition on ideals of the Noetherian ring $R$.
Consequently $f$ is injective as well as surjective, and is an isomorphism.
\end{proof}







